\section{结论要点}
\label{sec:takeaway}

把六天 1{,}648 个 commit 与 12 臂综合读数压成十条可复用结论。每条给出数据锚点与对后续工作的含义。

\begin{enumerate}
  \item \textbf{管线是净正的，不是成本项}。同一 80 格上，中文链+修复循环相对裸编译出 PDF +10 格、纯净 +21 格；compilebench 500 格中文臂纯净率反超裸编 +16.1pt。规范化+修复循环顺带修好了源文件自身的缺陷——「翻译引入的编译代价」被「规范化修复的源缺陷」覆盖后还有盈余。含义：对外口径可以直接说「过本管线的论文比原样编译更容易出 PDF」。

  \item \textbf{修复循环是单一最大杠杆，且收益可复现}。09-16 接线当日联合出 PDF +18.9pt（同口径复跑确认）；v2 对照 40 篇 34 个失败格救回 25（74\%）；本批真实臂 7 个管线失败格全数救回。143 条规则中装包/补件类贡献主要救回量。含义：规则沉淀速度（31$\to$143 条/5 天）就是出 PDF 率的爬坡速度，机制台账$\to$规则这条反馈环值得继续投入。

  \item \textbf{解析硬失败面已经枯竭，残面转为告警谱}。28{,}904 文件解析成功 100\%、一致率 99.99\%，唯一解析错误是语料源怪癖（tar 伪装成 .tex）。剩下的是告警级形态（\texttt{stray\_end}、\texttt{missing\_input} 各三千量级）与泄漏长尾（dollar 族 75 条）。含义：解析段从「修 bug」阶段进入「清长尾」阶段，优先级应让位于翻译与编译侧。

  \item \textbf{泄漏归零是样本框的函数，不是终态}。dollar 族在 1{,}955 文件核心层清零后，8 倍体量全量复测重现 76 条命中。教训：任何「清零」结论必须绑定样本框声明；核心层绿不代表长尾绿，全量复测节奏要保持。

  \item \textbf{翻译硬失败高度集中，提示词可定向修}。\texttt{cs\_dropped} 8 例中 7 例是 \verb|\ |、\verb|\,|、\verb|~| 三个脆弱间距命令被中文熔合吞掉——不是均匀散布的鲁棒性问题，是单一可点名机制。含义：在提示词里对这三个命令显式钉保护即可预期消灭该类大半；\texttt{ph\_miss} 的级联丢失（一格连丢 9 个占位符）提示要查译段级联失效而非逐点丢失。

  \item \textbf{质量缺陷呈结构性分布：术语是大头轻症、约定违例是小头重症}。术语族 263 条但 231 条轻症；「不可译」约定 25 条中 23 条重症——参考文献与命令名被误译是系统性协议违例。含义：术语表资产值得继续扩（轻症面大、边际收益高），约定违例则要靠校验规则硬拦（量小但每条都是产品级事故）。

  \item \textbf{「出 PDF $\neq$ 成功」有实证兜底}。Mode B 注入 132 处占位符破坏，编译前被校验链 132/132 全捕获。这是整条链路最关键的非对称保证：翻译错了宁可不出 PDF，不能出一份静默错位的 PDF。含义：校验层投入（检出率、误报率、延迟三指标全绿）是产品正确性的根基，不是开销。

  \item \textbf{评测器自证过稳定性}。计分卡 schema v3 重放 loop1 得 97.89/84.99 分毫不差；报告全链路可复算（落盘记录追加式 + 样本集过滤口径）。含义：跨日对比读数可信，「指标涨了」可归因到代码而非测量噪声——这是多日迭代能收敛的前提。

  \item \textbf{基线自己会动}。裸编译地板四天涨约 19pt（71.7$\to$90.2），原因是离线宏包资产与 tlmgr usertree 积累——工具链成熟本身抬高了所有臂的分母。含义：任何「相对基线的增益」都要标注基线时点；「裸引擎撑死七成」这类早期结论会过期，README 与文档里的静态断言要定期复测。

  \item \textbf{未达标项与短板是显式清单，不是模糊感受}。M2 纯净率差 1.25pt（88.75 vs 90）、死占位符 3 例、合平口径 93.7\%、alignbench 有效覆盖仅 42 对、tectonic 的 EPS/biber 双死穴、\texttt{cs\_dropped} 归因待定、计分卡过期行污染——七项各有明确下一步（归因、口径修正、路由前置、对照实验、工具参数）。含义：下一里程碑的验收条件已经写死，不存在「感觉差不多了」的空间。
\end{enumerate}
